\subsection{Rationalizing Denominators with Variables} 
Previously, we have rationalized denominators which contained roots of numbers.
\begin{example}
Simplify $$\sqrt{\dfrac{25}{12}}$$ as much as possible, including rationalizing the denominator.
\begin{lpic}{}{5}
\fractextleft{0.25}{-1}{$\sqrt{\dfrac{25}{12}}={}$};
\end{lpic}
\end{example}
\noi When working with variable expressions, it is easier to find a denominator whose root we can take.
\begin{example}
Simplify $$\dfrac{3x^3}{\sqrt[4]{y^9}}$$ by rationalizing the denominator. Assume all variables represent positive real numbers.
\begin{lpic}{}{3}
\fractextleft{0.25}{-2}{$\dfrac{3x^3}{\sqrt[4]{y^9}}={}$};
\end{lpic}
\end{example}
\noi If we factor coefficients into primes, we can handle them the same way that we do variables.
\begin{example}
Simplify $$\sqrt[3]{\dfrac{8x}{25y}}$$ as much as possible, including rationalizing the denominator. Assume all variables represent positive real numbers.
\begin{lpic}{}{5}
\fractextleft{0.25}{-4}{$\sqrt[3]{\dfrac{8x}{25y}}={}$};
\end{lpic}
\end{example}
\begin{note}
We can also ``rationalize the numerator'' using a similar procedure.
\end{note}